% Part: computability
% Chapter: recursive-functions
% Section: partial-functions

\documentclass[../../../include/open-logic-section]{subfiles}

\begin{document}

\olfileid[fr]{cmp}{rec}{par}
\olsection{Fonctions récursives partielles}

Pour motiver la définition des fonctions récursives, remarquons que
notre preuve de l'existence de fonctions calculables mais non
récursives primitives établit bien davantage. Elle n'utilisait
que la possibilité d'énumérer des fonctions \emph{totales}
$f_0,f_1,\dots$ de manière que $f_x(y)$ soit calculable en
fonction de $x$ et de~$y$. Le raisonnement s'applique donc à
\emph{toute classe de fonctions totales énumérable de cette
manière}. Nous voici devant une difficulté : nous aimerions
décrire explicitement toutes les fonctions calculables totales,
mais aucune énumération effective de telles fonctions ne peut
être exhaustive.\footnote{Précision éditoriale : la contradiction diagonale utilise la totalité de chaque fonction énumérée et une évaluation calculable uniforme. Elle ne contredit ni une description non effective des fonctions totales calculables ni l'énumération des fonctions partielles.}

Pour en sortir, il faut admettre les fonctions \emph{partielles}.
Nous verrons qu'il est possible d'énumérer les fonctions
partielles calculables.\iftag{TMs}{ En fait, nous le savons
presque déjà, puisque les machines de Turing peuvent être
énumérées systématiquement.}{} Nous reviendrons à l'argument
diagonal pour comprendre pourquoi il ne s'applique plus de
la même façon lorsque des fonctions partielles sont admises.

Que faut-il ajouter aux fonctions récursives primitives pour
obtenir toutes les fonctions récursives partielles ? Il faut
procéder en deux temps :
\begin{enumerate}
\item Adapter leur définition afin d'autoriser les fonctions
  partielles.
\item \emph{Ajouter} un schéma de définition qui produise
  de nouvelles fonctions partielles.
\end{enumerate}

La première étape est simple. Nous partons de nouveau de la
fonction nulle, du successeur et des projections, puis appliquons
la composition et la récurrence primitive. Il suffit d'adapter
ces deux opérations au cas où certains des termes en jeu ne
sont pas définis. Si $f$ et $g$ sont des fonctions partielles,
nous écrirons $f(x) \fdefined$ lorsque $f$ est définie en~$x$,
c'est-à-dire lorsque $x$ appartient à son domaine, et
$f(x) \fundefined$ dans le cas contraire. La notation
$f(x) \simeq g(x)$ signifie que les deux valeurs sont
indéfinies, ou qu'elles sont toutes deux définies et égales.
Nous emploierons ces notations également pour des termes plus
complexes. Si $h$ et $g_0,\dots,g_k$ sont des fonctions
partielles, nous convenons que
\[
h(g_0(\vec x),\dots,g_k(\vec x))
\]
est défini si et seulement si chaque $g_i$ est définie
en~$\vec x$ et si $h$ est définie au tuple de valeurs
$(g_0(\vec x),\dots,g_k(\vec x))$. Dans ces conditions, les
définitions de la composition et de la récurrence primitive
restent celles données plus haut, avec « $\simeq$ » à la
place de « $=$ ».

L'opération que nous ajoutons est la \emph{recherche non bornée},
ou \emph{minimisation non bornée}. Pour toute fonction partielle
$f(x,\vec z)$ à valeurs dans les entiers naturels, définissons
$\mu x \; f(x,\vec z)$ comme
\begin{quote}
  le plus petit $x$ tel que $f(0,\vec z),f(1,\vec z),
  \dots,f(x,\vec z)$ soient toutes définies et que
  $f(x,\vec z)=0$, s'il existe un tel~$x$.
\end{quote}
Dans le cas contraire, $\mu x \; f(x,\vec z)$ est indéfini.
Cette condition détermine la valeur de manière unique.

\begin{explain}
Notre définition ne fait appel ni\iftag{TMs}{ aux machines
de Turing ni}{} aux algorithmes, ni à un modèle de calcul
particulier. Comme pour la composition et la récurrence
primitive, la recherche non bornée possède néanmoins une
interprétation opératoire. Pour une fonction partielle
calculable, les arguments hors de son domaine correspondent
aux calculs qui ne s'arrêtent pas. Pour calculer
$\mu x \; f(x,\vec z)$, on calcule successivement
$f(0,\vec z),f(1,\vec z),f(2,\vec z),\dots$ jusqu'à obtenir~$0$.
Si l'un de ces calculs intermédiaires ne s'arrête pas,
la recherche ne s'arrête pas non plus.
\end{explain}

Pour toute relation $R(x,\vec z)$, posons
$\mu x \; R(x,\vec z)=\mu x \; (1 \tsub \Char{R}(x,\vec z))$.
Autrement dit, $\mu x \; R(x,\vec z)$ renvoie le plus petit~$x$
qui satisfait~$R$. Si $f(x,\vec z)$ est totale,
$\mu x \; f(x,\vec z)$ coïncide donc avec
$\mu x \; (f(x,\vec z)=0)$. La définition initiale est
cependant plus générale : $f$ peut n'être pas définie
partout, alors que la fonction caractéristique d'une
relation est toujours totale.

\begin{defn}
La classe des \emph{fonctions récursives partielles} est la plus
petite classe de fonctions partielles d'arités quelconques,
des entiers naturels dans les entiers naturels, qui contient
la fonction nulle, le successeur et les projections, et qui
est close par composition, récurrence primitive et recherche
non bornée.
\end{defn}

Certaines fonctions récursives partielles sont \emph{totales},
c'est-à-dire définies pour tous leurs arguments.

\begin{defn}
\ollabel{defn:recursive-fn}
La classe des \emph{fonctions récursives} est celle des
fonctions récursives partielles qui sont totales.
\end{defn}

On dit parfois « fonction récursive totale » pour souligner
qu'une fonction récursive est définie partout.

\end{document}
